\chapter[Differentiable Particle Filters]{Differentiable Particle Filters: Literature Survey and Claim Boundaries}
\label{ch:bf-dpf-literature-survey}

This chapter surveys the differentiable particle-filter literature relevant to
BayesFilter's program in
\texttt{docs/differentiable-particle-filter-program.md}.  The goal is not to
promote a single frontier method by enthusiasm alone.  The goal is to separate
what the literature supports, what the current student reports claim, what is
still only a project hypothesis, and what would be required before a
particle-based likelihood path becomes acceptable for HMC in nonlinear DSGE or
MacroFinance settings.

The survey therefore uses four evidence lanes:
\begin{enumerate}[label=(\roman*)]
  \item primary particle-filter and particle-flow literature;
  \item the CIP monograph chapters that already synthesize part of that
    literature in project notation;
  \item the two student baseline tracks now vendored under
    \texttt{experiments/student\_dpf\_baselines/vendor/};
  \item BayesFilter's current doctrinal constraint that target validity,
    derivative validity, and sampler validity are distinct questions.
\end{enumerate}
The student reports are useful because they expose implementation choices,
benchmark habits, and occasionally stronger or looser conclusions than the
primary literature.  They are not production authorities by themselves.  When
this chapter disagrees with them, it does so by naming the specific source or
contract distinction responsible for the difference.

\section{Scope and terminology}
\label{sec:bf-dpf-scope-terminology}

Write the hidden state as $x_t \in \mathbb{R}^{n_x}$, the observation as
$y_t \in \mathbb{R}^{n_y}$, and the structural parameter as
$\theta \in \mathbb{R}^{d_\theta}$.  The nonlinear state-space model is
\begin{align}
  x_t &\sim p_\theta(x_t \mid x_{t-1}), \\
  y_t &\sim p_\theta(y_t \mid x_t).
\end{align}
The filtering recursion propagates
\[
  p_\theta(x_t \mid y_{1:t-1})
  = \int p_\theta(x_t \mid x_{t-1})
    p_\theta(x_{t-1} \mid y_{1:t-1})\,dx_{t-1}
\]
and updates by Bayes' rule,
\[
  p_\theta(x_t \mid y_{1:t})
  \propto p_\theta(y_t \mid x_t)
  p_\theta(x_t \mid y_{1:t-1}).
\]
The marginal likelihood factorization is
\begin{equation}
\label{eq:bf-dpf-marginal-factorization}
  p_\theta(y_{1:T})
  = \prod_{t=1}^T p_\theta(y_t \mid y_{1:t-1}).
\end{equation}

A classical particle filter approximates the filtering law by a weighted atomic
measure,
\begin{equation}
\label{eq:bf-dpf-empirical-measure}
  \hat \pi_t^N(dx)
  = \sum_{i=1}^N w_t^{(i)}\,\delta_{x_t^{(i)}}(dx).
\end{equation}
Differentiable particle-filter proposals modify this basic construction in at
least three conceptually different ways:
\begin{enumerate}[label=(\alph*)]
  \item they change the \emph{proposal dynamics}, for example by transporting
    particles via a deterministic or stochastic flow;
  \item they change the \emph{resampling mechanism}, for example by replacing a
    categorical draw with a smooth relaxation or an optimal-transport map;
  \item they change the \emph{learning objective}, for example by optimizing a
    differentiable surrogate rather than maintaining an unbiased likelihood
    estimator.
\end{enumerate}
These distinctions are load-bearing.  A method can be differentiable without
being unbiased, useful for training without being suitable for HMC, or
value-consistent in a special case without being a generally target-preserving
likelihood path.

\section{Bootstrap particle filtering as the baseline}
\label{sec:bf-dpf-bootstrap-baseline}

The baseline reference remains the bootstrap particle filter of
\citet{gordon1993novel}, described at greater length in the CIP survey chapter
\texttt{ch26\_differentiable\_pf.tex}.  For particles
$\{x_t^{(i)}\}_{i=1}^N$, the basic recursion is:
\begin{enumerate}[label=(\roman*)]
  \item sample proposal particles
    $\tilde x_t^{(i)} \sim p_\theta(x_t \mid x_{t-1}^{(i)})$;
  \item assign weights
    $\tilde w_t^{(i)} \propto p_\theta(y_t \mid \tilde x_t^{(i)})$;
  \item normalize $w_t^{(i)} = \tilde w_t^{(i)} / \sum_j \tilde w_t^{(j)}$;
  \item resample ancestors according to $w_t^{(i)}$.
\end{enumerate}
The marginal-likelihood estimator is
\begin{equation}
\label{eq:bf-dpf-bpf-likelihood}
  \hat p_\theta^N(y_{1:T})
  = \prod_{t=1}^T \left( \frac{1}{N}\sum_{i=1}^N \tilde w_t^{(i)} \right).
\end{equation}
Under the standard bootstrap-particle-filter assumptions this estimator is
unbiased for the likelihood \citep{doucet2001sequential,andrieu2010particle}.
That statement is mathematically valuable because it gives BayesFilter a clear
reference distinction between:
\begin{itemize}
  \item an \emph{unbiased likelihood estimator},
  \item a differentiable surrogate for the same likelihood,
  \item and a value-only approximation with no exact-target claim.
\end{itemize}

The classical weakness is degeneracy.  As emphasized by
\citet{bengtsson2008curse} and \citet{snyder2008obstacles}, the effective
sample size can collapse rapidly in high dimension, especially when the
observation density is tight.  The student MLCOE report correctly highlights
this degeneracy problem in its literature review and uses it to motivate
particle-flow methods.  That part of the report is consistent with the standard
literature and with the CIP discussion.  The more ambitious conclusion one
cannot draw from degeneracy alone is that a differentiable replacement for
resampling automatically yields an HMC-suitable likelihood path.  Degeneracy
motivates algorithmic innovation; it does not settle the target-validity issue.

\section{Particle-flow filters}
\label{sec:bf-dpf-particle-flow-literature}

Particle-flow filters were introduced by \citet{daumhuang2008} and developed in
subsequent work by Daum and Huang to replace discrete resampling by a
continuous transport from a predictive density toward a posterior density.  The
CIP survey chapter presents the standard homotopy form
\begin{equation}
\label{eq:bf-dpf-homotopy}
  \log p_\lambda(x)
  = \log p(x \mid y_{1:t-1}) + \lambda \log p(y_t \mid x) + c(\lambda),
  \qquad \lambda \in [0,1],
\end{equation}
and then derives an affine flow field under Gaussian closure.  In the global
exact Daum--Huang (EDH) construction, the resulting ordinary differential
equation has the form
\begin{equation}
\label{eq:bf-dpf-edh-ode}
  \frac{dx}{d\lambda} = A(\lambda)x + b(\lambda),
\end{equation}
with coefficients determined by the current predictive covariance and the local
likelihood precision.

The literature point that matters for BayesFilter is not merely the formula but
the \emph{status} of the formula.  Under linear-Gaussian assumptions the flow
recovers the Kalman update, so EDH supplies an exact special-case benchmark.
Outside that setting, the global Gaussian closure is an approximation.  The CIP
chapter states this carefully, and BayesFilter should preserve that wording.
The MLCOE report's literature review and implementation framing are broadly
consistent with this distinction: it treats EDH as analytically appealing but
still motivates localized or corrected variants for non-Gaussian filtering.

\subsection{Local EDH and particle-specific linearization}
\label{sec:bf-dpf-ledh-literature}

\citet{li2017particle} and related particle-flow literature replace the global
linearization with particle-specific local linearizations, producing what is
commonly called local EDH (LEDH).  The CIP methodology chapter writes the local
homotopy coefficients in terms of a particle-specific Jacobian and local
precision matrix.  This reduces the mismatch between the true nonlinear
likelihood geometry and the global affine closure, but it changes the algorithm
from a single deterministic transport field into a cloud of particle-specific
transports.

The main literature-supported benefit is improved adaptability to nonlinear
observation structure.  The main literature-supported cost is that once the flow
is discretized and localized, it is no longer generally exact.  The student
repositories both pay serious attention to LEDH.  The MLCOE report uses LEDH as
a core component of its PF-PF implementation track; the
\texttt{advanced\_particle\_filter} repository also treats EDH/LEDH as a major
classical and differentiable baseline.  Where their conclusions become stronger
than the primary literature is usually not in claiming exactness of the local
flow, but in how quickly they move from improved filtering behavior to stronger
HMC conclusions.  BayesFilter should keep that boundary explicit.

\subsection{PF-PF importance correction}
\label{sec:bf-dpf-pfpf-literature}

The particle-flow particle filter (PF-PF) literature, especially
\citet{li2017particle}, treats the deterministic flow not as the final target
itself but as a proposal mechanism inside an importance-sampling scheme.  If a
flow map $\Phi$ transforms a proposal particle $x_{t,0}$ into a transported
particle $x_{t,1}=\Phi(x_{t,0})$, then the corrected weight is obtained by the
usual change-of-variables formula:
\begin{equation}
\label{eq:bf-dpf-pfpf-weight}
  \tilde w_t
  \propto
  \frac{p_\theta(y_t \mid x_{t,1})\,p_\theta(x_{t,1} \mid x_{t-1})}
       {q_\Phi(x_{t,1} \mid x_{t-1},y_t)}
  =
  \frac{p_\theta(y_t \mid x_{t,1})\,p_\theta(x_{t,1} \mid x_{t-1})
        \left|\det D\Phi(x_{t,0})\right|}
       {p_\theta(x_{t,0} \mid x_{t-1})}.
\end{equation}
This is the literature step that makes BayesFilter's agreed development ladder
sensible.  The flow alone is not enough; the importance correction is the first
step that tries to recover a principled proposal/target relationship.

The student reports are useful here because they highlight two different
implementation emphases.  The MLCOE report foregrounds PF-PF consistency and
Jacobian correction as a central motivation.  The
\texttt{advanced\_particle\_filter} codebase, especially its classical thread,
also keeps EDH/LEDH and PF-PF in view, but its newer TensorFlow notebooks shift
more attention toward differentiable-resampling and HMC experiments.  This is a
useful warning for BayesFilter: once differentiable objectives become available,
it becomes easy to underemphasize the correction logic that preserves a clear
proposal interpretation.

\section{Kernel and stochastic particle-flow work}
\label{sec:bf-dpf-kernel-stochastic-literature}

The student MLCOE report covers two additional literatures that are helpful even
when they are not BayesFilter's primary intended production path.

First, kernel-embedded particle-flow methods such as the work emphasized by
\citet{hu2021particle} address high-dimensional or sparsely observed systems by
constructing flow fields in a reproducing-kernel Hilbert space.  The report's
treatment of marginal collapse in standard scalar-kernel formulations and the
move to matrix-valued kernels is useful as a reminder that transport ideas can
fail in underobserved subspaces even when they appear numerically stable on
observed coordinates.

Second, stochastic particle-flow work, including Dai and Daum's stiffness and
control-inspired homotopy-schedule analyses, is useful for BayesFilter because
it names a real implementation hazard: the flow ODE may become numerically
stiff when observation precision is high.  The CIP chapters and the student
reports are aligned on this issue.  This does not yet make stochastic
particle-flow the preferred BayesFilter path, but it justifies carrying
stiffness as a first-class diagnostic obligation in any EDH or LEDH
implementation.

\section{Differentiable resampling literature}
\label{sec:bf-dpf-differentiable-resampling-literature}

The differentiable-particle-filter frontier is driven by the resampling
bottleneck.  Standard multinomial or systematic resampling is a discrete map,
so a naive pathwise derivative with respect to the weights is zero almost
everywhere.  The literature response has taken at least three main forms.

\subsection{Soft resampling and related smooth relaxations}
\label{sec:bf-dpf-soft-resampling}

One family of methods uses a soft or partially collapsed replacement for hard
ancestor selection.  The CIP resampling chapter follows
\citet{zhumurphyjonschkowski2020} and writes a representative soft-resampling
update as
\begin{equation}
\label{eq:bf-dpf-soft-resampling}
  \tilde x^{(i)}
  = \alpha x^{(a_i)} + (1-\alpha)\bar x_w,
  \qquad
  \bar x_w = \sum_{j=1}^N w^{(j)}x^{(j)}.
\end{equation}
The mean of this soft-resampled point agrees with the weighted empirical mean,
but nonlinear functionals of the cloud are biased unless $\alpha=1$ and one
reverts to standard resampling.  This is precisely the kind of statement that
BayesFilter should preserve because it names the central trade-off: smooth
gradients are purchased by changing the target.

The \texttt{advanced\_particle\_filter} repository makes this trade-off very
explicit in its HMC-with-differentiable-resampling thread.  Its notebook index
states that the linear-Gaussian-observation diagnostics notebook is used to
compare DPF gradient quality against Kalman ground truth, specifically to
isolate DPF bias from observation nonlinearity.  This is a strong design move,
and BayesFilter should reuse the idea.  The important doctrinal qualification is
that such a diagnostic identifies and measures bias; it does not erase it.

\subsection{Entropy-regularized optimal transport}
\label{sec:bf-dpf-ot-resampling}

A second family replaces categorical resampling by an optimal-transport map.
The central reference for BayesFilter is \citet{corenflos2021differentiable},
which formulates resampling as a transport problem between a weighted empirical
measure and a uniformly weighted empirical measure.  The CIP chapter writes the
entropically regularized problem as
\begin{equation}
\label{eq:bf-dpf-entropic-ot}
  \Pi_\varepsilon^*
  = \arg\min_{\Pi \in \mathcal U(w,u)}
    \langle \Pi, C \rangle - \varepsilon H(\Pi),
\end{equation}
where $w$ is the source weight vector, $u$ is the uniform target weight vector,
$C$ is the transport-cost matrix, and $H(\Pi)$ is the entropy term.  The
barycentric projection of the resulting coupling defines a differentiable
replacement for resampling.

The literature-supported gain is a smooth transport map with a clear geometric
interpretation.  The literature-supported cost is that the regularized
transport is not the same object as exact resampling; it is a relaxed
projection whose bias depends on the regularization and on the finite Sinkhorn
solve.  The \texttt{advanced\_particle\_filter} repository is especially useful
here because it records an explicit parameter sweep over the entropy parameter
and Sinkhorn iterations, and its amortized-operator note states approximate
speedups and residual approximation error.  This helps BayesFilter see where
engineering enthusiasm can outrun target-discipline: faster wall clock is not
yet a target-preservation theorem.

\subsection{Amortized or learned OT operators}
\label{sec:bf-dpf-amortized-ot}

A third family learns to approximate the OT resampling map directly.  The
\texttt{advanced\_particle\_filter} repository's
\texttt{docs/amortized\_ot\_operator.md} is a particularly useful internal
comparison object because it is unusually explicit about the architecture,
training distribution, wall-clock speedups, and approximation residuals.  Its
claimed use is to replace iterative Sinkhorn solves by a learned set-based
operator that takes weighted clouds and returns approximately barycentric
uniform clouds.

This line of work is relevant to BayesFilter because it illustrates the final
rung of the agreed ladder: neural OT as an acceleration layer on top of OT
resampling, not as the first production baseline.  It also illustrates the key
reason BayesFilter remains cautious.  Once the resampling operator is learned,
any posterior change induced by operator approximation must be counted as part
of the target definition or explicitly corrected.  The advanced student repo is
helpfully honest about this: its architecture note reports a nonzero
approximation error and a visible posterior-mean shift relative to Sinkhorn and
Kalman comparisons.  That honesty is useful, and BayesFilter should preserve it
rather than smoothing it away in prose.

\section{Learning-based differentiable particle filters}
\label{sec:bf-dpf-learning-based-literature}

The student MLCOE report places differentiable particle filters in a broader
learning-based literature that includes \citet{jonschkowski2018differentiable},
\citet{karkus2018particle}, and related end-to-end filtering architectures.
Those works are useful because they clarify a distinct goal: some DPF methods
are designed primarily to train transition, observation, or proposal networks
within algorithmic priors.  In that regime the objective is often a
training-friendly surrogate, not an exact Monte Carlo likelihood estimator.

This difference matters for BayesFilter's HMC program.  A learning objective can
be perfectly respectable for representation learning and still be the wrong
object for Hamiltonian Monte Carlo.  The MLCOE report groups differentiable and
learning-based filtering together in a broad frontier category, which is fine as
a literature map.  BayesFilter needs a sharper split: methods designed for
end-to-end learning should not be treated as HMC-ready merely because they
provide gradients.

\section{Particle methods for parameter inference}
\label{sec:bf-dpf-parameter-inference-literature}

Parameter inference with particle methods occupies a wider space than the
BayesFilter DPF program.  The classical pseudo-marginal reference is
\citet{andrieu2010particle}, which uses an unbiased particle-filter likelihood
estimator inside a Metropolis--Hastings correction.  The student MLCOE report
also discusses particle marginal Metropolis--Hastings and particle HMC as bonus
questions, and the \texttt{advanced\_particle\_filter} repository devotes a full
thread to HMC with differentiable resampling.

The survey conclusion BayesFilter should draw is careful but not negative:
\begin{enumerate}[label=(\roman*)]
  \item pseudo-marginal particle MCMC has a clear exact-target story when its
    assumptions hold;
  \item differentiable particle filters may be useful for gradient-based
    parameter inference, but their exact target status depends on how the flow,
    resampling, and learning approximations are handled;
  \item the current literature does not justify collapsing all of these methods
    into a single statement such as ``particle HMC is correct whenever the
    filter is differentiable.''
\end{enumerate}
This is the main place where BayesFilter's doctrine remains intentionally more
conservative than the tone of some experimental reports.

\section{Comparison against the student literature reviews}
\label{sec:bf-dpf-student-comparison}

The two student tracks are genuinely useful as comparison objects because they
cover partially different slices of the literature.

The MLCOE report contributes a broad survey arc from Kalman and UKF through SMC,
EDH, LEDH, PF-PF, kernel particle flow, stochastic particle flow,
differentiable particle filters, and particle MCMC.  Relative to BayesFilter's
current monograph, it usefully foregrounds:
\begin{itemize}
  \item kernel-embedded particle flow and marginal-collapse concerns;
  \item stiffness mitigation in particle-flow ODEs;
  \item a larger spread of geophysical benchmark literature.
\end{itemize}
Its weakest points are not the range of topics but the tendency to drift from
survey statements into architecture or application enthusiasm without always
maintaining BayesFilter's distinction between filtering quality, target status,
and HMC validity.

The \texttt{advanced\_particle\_filter} repository contributes a more explicit
engineering and experimental arc.  Relative to both CIP and the MLCOE report,
it is especially helpful about:
\begin{itemize}
  \item linear-Gaussian diagnostics that isolate DPF gradient bias;
  \item Sinkhorn-parameter sweeps;
  \item MAP, Laplace, and HMC pipeline experiments;
  \item explicit amortized-OT approximation costs and speedups.
\end{itemize}
Its main limitation for BayesFilter is that it remains a research codebase and
not a production contract.  Many conclusions are notebook- and benchmark-based,
which makes them valuable hypotheses and diagnostics but not enough by
themselves to upgrade BayesFilter's HMC claims.

On balance, the student materials and the primary literature agree on the
following broad story: particle-flow and differentiable-resampling methods are a
serious research path for nonlinear filtering and gradient-based parameter
inference, especially where Gaussian closures are inadequate.  Where BayesFilter
remains more conservative is in the final interpretive step.  The literature and
the student reports motivate the path; they do not yet eliminate the need for
per-rung target, gradient, variance, and compiled-path audits.

\section{Method-family claim register}
\label{sec:bf-dpf-claim-register}

Table~\ref{tab:bf-dpf-claim-register} summarizes the method families most
relevant to the BayesFilter program.

\begin{table}[ht]
\centering
\small
\begin{tabular}{p{0.18\textwidth} p{0.18\textwidth} p{0.16\textwidth} p{0.34\textwidth}}
\hline
Method family & Likelihood / target status & Gradient status & Main BayesFilter blocker \\
\hline
Bootstrap PF & Unbiased likelihood estimator under standard assumptions & Not pathwise differentiable through standard resampling & Degeneracy and nondifferentiable resampling \\
EDH flow alone & Exact only in linear-Gaussian recovery; otherwise controlled approximation & Smooth value path if implemented numerically stably & No generic corrected-target claim outside special cases \\
LEDH flow alone & Improved local approximation, not generic exactness & Smooth local-flow path & Local linearization and discretization bias \\
PF-PF with flow weights & Proposal-correction path with principled importance weighting & Depends on flow/Jacobian implementation & Variance, Jacobian accuracy, and practical stability \\
Soft resampling & Biased differentiable surrogate & Pathwise-friendly & Target bias must be named explicitly \\
Entropic OT resampling & Biased differentiable transport relaxation & Differentiable through Sinkhorn / OT solver & Relaxation bias and solver approximation \\
Neural OT resampling & Learned surrogate to a relaxed OT map & Differentiable neural operator & Additional approximation layer and training-distribution mismatch \\
\hline
\end{tabular}
\caption{Claim register for the main differentiable particle-filter method families discussed in this chapter.}
\label{tab:bf-dpf-claim-register}
\end{table}

\section{Implication for the BayesFilter program}
\label{sec:bf-dpf-literature-implication}

The literature and student comparisons support the development ladder already
recorded in \texttt{docs/differentiable-particle-filter-program.md}: start with
EDH linear-Gaussian recovery, add PF-PF importance correction, then study soft
or OT resampling, and only later consider LEDH plus neural OT.  This ordering is
not arbitrary.  It matches the places where the literature moves from a clearer
proposal or benchmark interpretation toward increasingly relaxed,
learning-dependent surrogates.

The survey therefore supports the BayesFilter path as a \emph{research and
engineering program}.  It does not yet support skipping the intermediate gates.
In particular, none of the surveyed material removes the need to test, rung by
rung, whether the scalar being differentiated is an exact likelihood, an
unbiased estimator, a biased surrogate, or merely a useful engineering score.
